problem:  
Let \( M^{2n} \) be a compact, oriented, even-dimensional Riemannian manifold with strictly positive sectional curvature. Assume that \( M \) admits an effective isometric \( S^1 \)-action whose fixed-point set is finite.  

**Conjecture:** Among all such manifolds, the Euler characteristic \( \chi(M) \) satisfies  
\[
\chi(M) \leq n+1,
\]
with equality if and only if \( M \) is diffeomorphic to the complex projective space \( \mathbb{C}P^n \) equipped with a linear \( S^1 \)-action.  

Why is it a "good" problem:  
This problem proposes a quantitative refinement of the Hopf conjecture within the framework of positively curved manifolds with symmetry. Hopf’s conjecture states that even-dimensional, positively curved manifolds must have positive Euler characteristic—an assertion of sign. The present problem extends this by seeking an *upper bound* on the Euler characteristic in the presence of an isometric \( S^1 \)-action, thereby exploring whether symmetry and positive curvature together impose strong topological rigidity.  

The question sits at the nexus of several central domains:  
- **Riemannian geometry:** through the constraint of positive sectional curvature,  
- **Transformation group theory:** via isometric \( S^1 \)-actions,  
- **Equivariant topology and fixed-point theory:** since \( \chi(M) \) is determined by the local data of the fixed points, and  
- **Topological classification:** asking whether \( \mathbb{C}P^n \) emerges as the extremal or “maximally symmetric” case.  

The conjecture echoes the pattern known for homogeneous and biquotient examples—all known positively curved manifolds with circle actions have Euler characteristic at most \( n+1 \)—but no general proof is known. Proving or disproving this would sharpen our understanding of how curvature restricts topology under group actions, bridging global Riemannian geometry with equivariant index theory.  

Its formulation is simple (“Can we bound \( \chi(M) \) by \( n+1 \)?”), yet its resolution would require deep synthesis across geometry, topology, and representation theory—hallmarks of a genuinely profound and forward-looking mathematical problem.